%% ---------------------------------------------------------------------------
%% Revised Section 6.
%% ---------------------------------------------------------------------------
\section{Conclusion and Future Work}

We asked three questions of a protocol that fixes the decision function (LLM-as-teacher), runs every explainer on every representation under a common budget, sparsity and receptive field, and evaluates every detector against majority-class and confidence baselines. To the first, whether the representation or the explainer determines the signature of an explanation, the answer is the explainer: within the budget every explanation is sufficient and almost none is necessary on every representation, and the degrees differ between attribution methods, not between discrete and continuous representations. To the second, whether explanation features detect errors beyond the model's confidence, the answer is no for the language model in any configuration or combination, for either definition of error, with linear or non-linear detectors and among its over-confident errors; no graph topology or explainer outperforms the language model's own explainer. To the third, what a distilled surrogate inherits from its teacher, the answer is the decisions and confidence in them, not the errors, which a gold-trained twin makes too, nor the reasoning, which its attributions do not reflect beyond chance. The only place where explanation features help is this surrogate, whose confidence hides the inherited errors, and the help never carries it to the level of its teacher.

These results do not contradict the gains of explanation-based calibration under domain shift~\cite{ye2022calibrating}, where the confidence is less reliable and the calibrator is trained on target-domain examples; on in-distribution data, where confidence is informative, the explanation metrics we study are largely redundant with it.

\paragraph{\textbf{Scope and limitations}}
The design is controlled (2 explainers \(\times\) 4 topologies \(\times\) 2 datasets plus the language model, a node-feature ablation and three controls with over fifty additional surrogates) and the pipeline is released (\url{https://github.com/FabioYanezRomero/explanation-signatures-error-detection}), so the protocol can serve as a testbed for explanation-based error detection; we recommend that such studies report majority-class and confidence baselines, state the class that is explained, and stratify only by quantities available at inference time. The limitations are those of the setting: two English benchmarks, one language model (BERT) and one GNN architecture (GCN), a 20\% sparsity budget, and Shapley-style explainers only. A stronger surrogate, a different sparsity, or gradient-based explainers could change the descriptive signatures; we do not expect them to change the error-detection result, which is driven by the redundancy between explanation metrics and confidence, not by any property of a particular explainer.

\paragraph{\textbf{Future work}}
Three directions follow. First, explanation features may complement confidence in settings where confidence is known to fail, such as distribution shift, where explanation-based calibrators have been reported to help~\cite{ye2022calibrating}, or adversarial inputs, which our benchmarks do not cover. Second, the gain of explanation features over a surrogate's confidence among its over-confident errors on AG News deserves a dedicated study with more errors per regime, in particular whether it survives when the surrogate is calibrated against gold labels. Third, the over-confidence that distillation produces suggests that student models should be calibrated against gold labels, not only against their teacher, before being used as proxies for the teacher's behaviour.
